\section{Privacy Impact}
MANTA's privacy role follows from its threat-detection role. The detector needs telemetry to find
suspicious behavior, but that telemetry can itself reveal user and application behavior. The system
therefore enforces a metadata-only telemetry boundary: packet payloads are not stored or exported,
and the released views progressively suppress direct identifiers and strong destination hints. From
a GDPR-oriented perspective \citep{gdpr_2016_679}, this aligns with data minimization, purpose
limitation, and storage-limitation principles. The design still processes behaviorally revealing
metadata, so transparency, consent, and retention control remain necessary.

MANTA implements several telemetry-reduction mechanisms:
\begin{itemize}[leftmargin=*]
  \item metadata-only capture instead of payload inspection,
  \item hashed or removed direct identifiers in the reduced views,
  \item coarse or distribution-aware bucketization of many released features,
  \item privacy-student learning over a reduced feature view,
  \item federated-style learning over privacy-student latents instead of raw full-feature tables,
  \item explicit auditing of both release privacy and observer leakage.
\end{itemize}

These mechanisms reduce what MANTA releases or shares. They are not equivalent to a formal privacy
guarantee.

The evaluation also distinguishes hashing from deletion. Hashing removes direct readability of a
field, but deterministic salted hashes can still preserve stable equality and fingerprint
structure. Deleting a suppressed feature removes it from the released information set entirely. The
hash-versus-delete experiment in \Cref{tab:hash-delete} therefore treats hashing as
pseudonymization rather than anonymization. Hashing can recover detection utility, but it may
preserve learnable app-family or app-ID signal for an attacker.

\section{Ethical Considerations}
Responsible deployment requires:
\begin{itemize}[leftmargin=*]
  \item informed consent before monitoring starts,
  \item understandable disclosure of what is and is not collected,
  \item limited retention with user-triggered purge capability,
  \item strict access controls on backend triage endpoints,
  \item constraints against repurposing telemetry for non-security surveillance.
\end{itemize}
These points are consistent with EDPB guidance on transparency and lawful processing in digital
services \citep{edpb_guidelines_2020}. Security hardening practices from mobile application
standards, such as OWASP MASVS, further reduce secondary risk \citep{owasp_masvs}.

\section{Release Privacy versus Observer Leakage}
The release and observer measurements expose different failure modes. In the release view,
\texttt{medium} and \texttt{strict} substantially reduce exact app re-identification. Those same views
still leak high-level semantics such as app family and destination behavior. The observer benchmark
then shows that a passive eavesdropper can recover strong information from encrypted-flow sequences
even when the exported representation is reduced.

A release-privacy defender controls the representation that is exported or shared. A passive
observer does not consume that exported representation; they consume the original traffic behavior
on the network path. This is why export-time coarsening does not neutralize traffic-analysis
attacks. Works such as Deep Fingerprinting and later evaluations of website-fingerprinting defenses
reinforce this distinction \citep{deep_fingerprinting_2018,wf_sok_2022}.

\section{Interpretation of the Negative Result}
Feature-level coarsening is not enough to provide strong privacy against passive observers. Many
metadata-only monitoring designs rely on payload omission as their primary privacy argument. The
MANTA results show that this argument is incomplete. The released representation can be made less
revealing, but encrypted traffic remains fingerprintable at the observer level. This is a limitation
result, not a claim that MANTA solves passive traffic analysis.

\section{Leakage and Utility Evaluation}
MANTA benchmarks privacy risk with two complementary views:
\begin{itemize}[leftmargin=*]
  \item utility and leakage under feature-reduced release views, and
  \item observer leakage from encrypted-flow sequence fingerprinting.
\end{itemize}
The two measurements keep utility and inference risk visible at the same time. A reduced export
format has to retain threat-detection signal while reducing what secondary attackers can infer.

\section{Limitations and Threats to Validity}
\begin{itemize}[leftmargin=*]
  \item The anomaly models improve substantially in the public-corpus backend run, but the result is
  not a full adaptive-window evaluation of every model family.
  \item The Android-deployable models use an adaptive-window protocol, while the backend and privacy
  reports use practical source-balanced caps. These results should be compared within their
  documented protocols rather than treated as one fully unified full-corpus run.
  \item Full adaptive backend training was attempted but exceeded practical memory/time on the
  6.0M-flow corpus. The completed backend/privacy reports therefore use documented
  source-balanced raw-flow caps.
  \item The Adaptive hybrid RF runtime hybrid is evaluated on a bounded adaptive-window holdout, but
  still needs faster full-corpus scoring and a dedicated on-device alert-quality study for the final
  alert-evidence policy.
  \item The dataset mix is public and heterogeneous, but it still inherits each source dataset's
  collection biases and labeling noise.
  \item The privacy-student and federated paths are measured prototypes, not production defenses.
  \item Android foreground-service type declarations and Play Console review requirements create an
  additional deployment constraint for VPN-style monitoring on recent Android versions.
  \item No packet padding, dummy traffic, or timing obfuscation is implemented, so passive-observer
  protection remains out of scope.
  \item Additional privacy attacks such as membership inference, poisoning, and adaptive
  counterfactual adversaries remain future work.
\end{itemize}
